% Zone „Das kennst du schon“ – aus bank/winkel-dreiecke/zone.jsonl (zusammenbau v0.1)
\einheitenkopf[zone][Kennst du schon]{Das kennst du schon}
\setcounter{aufgabe}{0}

%% TODO Titel: Ich-kann-Satz fehlt in der Bank (Platzhalter: Kettenname) – Nr. 1
%% TODO Anweisung: ein Satz über der ersten Teilaufgabe fehlt in der Bank – Nr. 1
\begin{aufgabe}{Mit Lineal und Geodreieck zeichnen}
\begin{teile}
\teil Miss mit dem Lineal die Länge und die Breite des Rechtecks.

\rechteck{5}{3}
Länge: \leerfeld[cm] , Breite: \leerfeld[cm]
\teil Zeichne eine Gerade g und einen Punkt P, der 3 cm von g entfernt liegt. Zeichne mit dem Geodreieck die Senkrechte zu g durch P und die Parallele zu g durch P.

\begin{ksys}[xmin=0,xmax=10,ymin=0,ymax=6] \end{ksys}
\end{teile}
\end{aufgabe}

%% TODO Titel: Ich-kann-Satz fehlt in der Bank (Platzhalter: Kettenname) – Nr. 2
%% TODO Anweisung: ein Satz über der ersten Teilaufgabe fehlt in der Bank – Nr. 2
\begin{aufgabe}{Addieren und Subtrahieren im Bereich bis zum Vollwinkel, auch mit Dezimalzahlen (Vielfaches minus zwei Dezimalgrade)}
\begin{teile}
\teil $38^\circ + 29^\circ$ \leerfeld °
\teil $96{,}4^\circ + 38{,}7^\circ$ \leerfeld °
\end{teile}
\end{aufgabe}

%% TODO Titel: Ich-kann-Satz fehlt in der Bank (Platzhalter: Kettenname) – Nr. 3
%% TODO Anweisung: ein Satz über der ersten Teilaufgabe fehlt in der Bank – Nr. 3
\begin{aufgabe}{Längeneinheiten m und km umrechnen (1,5 km = 1500 m)}
\begin{teile}
\teil 3 km in Meter \leerfeld[m]
\teil 2,4 km in Meter \leerfeld[m]
\end{teile}
\end{aufgabe}

%% TODO Titel: Ich-kann-Satz fehlt in der Bank (Platzhalter: Kettenname) – Nr. 4
%% TODO Anweisung: ein Satz über der ersten Teilaufgabe fehlt in der Bank – Nr. 4
\begin{aufgabe}{Vierecksarten kennen}
\begin{teile}
\teil Wie heißt ein Viereck mit vier rechten Winkeln und vier gleich langen Seiten? \leerfeld
\teil Welche Seiten des Vierecks sind parallel? Wie heißt das Viereck?

\trapez{5}{3}{2.5}
\end{teile}
\end{aufgabe}

%% TODO Titel: Ich-kann-Satz fehlt in der Bank (Platzhalter: Kettenname) – Nr. 5
%% TODO Anweisung: ein Satz über der ersten Teilaufgabe fehlt in der Bank – Nr. 5
\begin{aufgabe}{Kreis mit dem Zirkel zeichnen, Radius und Durchmesser unterscheiden}
\begin{teile}
\teil Zeichne mit dem Zirkel einen Kreis mit dem Radius 3 cm.

\begin{ksys}[xmin=0,xmax=10,ymin=0,ymax=6] \end{ksys}
\teil Welche Strecke ist der Radius r, welche der Durchmesser d? Miss beide.

\begin{kreis}[2] \mittelpunkt{M} \radius{40}{r} \durchmesser{0}{d} \end{kreis}
r = \leerfeld[cm] , d = \leerfeld[cm]
\end{teile}
\end{aufgabe}

%% TODO Titel: Ich-kann-Satz fehlt in der Bank (Platzhalter: Kettenname) – Nr. 6
%% TODO Anweisung: ein Satz über der ersten Teilaufgabe fehlt in der Bank – Nr. 6
\begin{aufgabe}{Addieren und Subtrahieren im Bereich bis zum Vollwinkel, auch mit Dezimalzahlen (Vielfaches minus zwei Dezimalgrade) – Fehler finden}
\begin{teile}
\teil Mia rechnet: \rechnung{180^\circ - 46{,}8^\circ - 37{,}5^\circ &= 96{,}3^\circ} Finde den Fehler und rechne richtig.
\end{teile}
\end{aufgabe}

%% TODO Titel: Ich-kann-Satz fehlt in der Bank (Platzhalter: Kettenname) – Nr. 7
%% TODO Anweisung: ein Satz über der ersten Teilaufgabe fehlt in der Bank – Nr. 7
\begin{aufgabe}{Addieren und Subtrahieren im Bereich bis zum Vollwinkel, auch mit Dezimalzahlen (Vielfaches minus zwei Dezimalgrade) – selbst rechnen}
\begin{teile}
\teil $180^\circ - 38{,}6^\circ - 55{,}7^\circ$ \leerfeld °
\end{teile}
\end{aufgabe}
